\documentclass[../../main.tex]{subfiles}
\begin{document}
\section{Mayer--Vietoris 定理}\label{section:ch17-the-mayer-vietoris-theorem}

我们现在叙述一个一般定理，它把一个流形的 de Rham 群用其两个开子集的 de Rham 群表示出来。

设 $M$ 是一个光滑流形，$U,V\subseteq M$ 是两个开子集，满足 $M=U\cup V$。考虑包含映射
\begin{align*}
\begin{array}{ccccc}
&& M &&\\
& \overset{j}{\nearrow} && \overset{k}{\searrow}\\
U &&&& V\\
& \underset{i}{\searrow} && \underset{l}{\nearrow}\\
&& U\cap V &&
\end{array}
\end{align*}
这些包含映射在 de Rham 上同调上诱导出映射。

\begin{theorem}[Mayer--Vietoris 定理]\label{theorem:ch17-mayer-vietoris}
设 $M$ 是一个光滑流形，$U,V$ 是 $M$ 的两个开子集且 $M=U\cup V$。则对每个 $p$，存在一个连接同态 $\delta^*:H^p_{\mathrm{dR}}(U\cap V)\to H^{p+1}_{\mathrm{dR}}(M)$，使得下面的序列正合：
\begin{align}
\cdots\to H^p_{\mathrm{dR}}(M)\xrightarrow{k^*\oplus l^*}H^p_{\mathrm{dR}}(U)\oplus H^p_{\mathrm{dR}}(V)
\xrightarrow{i^*-j^*}H^p_{\mathrm{dR}}(U\cap V)\xrightarrow{\delta^*}H^{p+1}_{\mathrm{dR}}(M)\to\cdots,
\end{align}
其中 $i^*-j^*$ 把 $(\omega,\eta)$ 映为 $i^*\omega-j^*\eta$。
\end{theorem}

\begin{figure}[ht]
\centering
\includegraphics[width=0.6\textwidth]{fig-17-2.png}
\caption{Mayer--Vietoris 覆盖}\label{figure:ch17-mayer-vietoris-cover}
\end{figure}

\begin{proof}
证明在下一节给出。基本想法是：先在形式层面构造一个短正合列
\begin{align*}
0\to\Omega^p(M)\to\Omega^p(U)\oplus\Omega^p(V)\to\Omega^p(U\cap V)\to0,
\end{align*}
再用蛇引理（zigzag lemma）导出上同调的长正合列。
\end{proof}

\begin{example}[球面的 de Rham 群]\label{example:ch17-sphere-cohomology}
利用 Mayer--Vietoris 定理，可以归纳地算出球面的 de Rham 群：
\begin{align*}
H^0_{\mathrm{dR}}(S^n)&\cong\mathbb{R},\\
H^n_{\mathrm{dR}}(S^n)&\cong\mathbb{R}\quad(n\geqslant1),\\
H^p_{\mathrm{dR}}(S^n)&=0\quad(p\neq0,n).
\end{align*}
把 $S^n$ 写成两个略形变的开半球之并，它们的交微分同胚于 $S^{n-1}\times\mathbb{R}$，从而与 $S^{n-1}$ 同伦。对 Mayer--Vietoris 序列作归纳即得上式。
\end{example}

\begin{figure}[ht]
\centering
\includegraphics[width=0.6\textwidth]{fig-17-3.png}
\caption{球面的 Mayer--Vietoris 覆盖}\label{figure:ch17-sphere-cover}
\end{figure}

\begin{proposition}[第一次 de Rham 群与基本群]\label{proposition:ch17-h1-fundamental-group}
设 $M$ 是一个连通光滑流形，$q\in M$。存在一个自然映射从 $H^1_{\mathrm{dR}}(M)$ 到 $\operatorname{Hom}(\pi_1(M,q),\mathbb{R})$，它把一个闭 $1$-形式的上同调类映为“沿环路积分”给出的同态。
\end{proposition}
\end{document}
